%%
% 第二章 相关工作与问题定义
% 综述章：只引用真实文献，不做 TODO 占位
%%

\chapter{相关工作与问题定义}
\label{cha:relatedwork}

本章综述酶底物特异性预测的相关工作，介绍本研究所采用的 EZSpecificity 参考模型框架，整理其中所涉及的蛋白、分子与几何表征学习组件，并在综述基础上精确定义本研究的问题与总体方案。综述范围限于本论文实验所直接涉及的方法层内容，不作全领域大综述；未在本研究中使用的技术路线（如反应级预测、区域选择性建模等）仅在第六章未来工作中提及。

\section{酶与底物特异性预测方法综述}
\label{sec:methods-survey}

酶底物特异性预测的计算方法可按信息源与任务粒度分为若干条相互关联的技术路线。从信息源看，方法大致划分为基于序列、基于分子图、基于三维结构与三者混合四类；从任务粒度看，可分为酶功能分类（如 EC 编号预测）、酶--底物二分类、酶--底物回归（如 $k_{\mathrm{cat}}$ 预测）与反应级产物预测。以下择取与本研究直接相关的代表性工作综述。

\textbf{基于序列的方法与蛋白语言模型}。基于蛋白序列的方法以 BLAST 同源搜索\cite{BLAST1990} 为传统代表，近年则由蛋白语言模型主导。Ryu 等提出的 DeepEC\cite{DeepEC2019} 以卷积神经网络对酶的氨基酸序列编码，实现对酶学委员会编号（Enzyme Commission number，EC 编号）的高通量预测，在代谢组学与合成生物学筛选中已成为常用工具；Yu 等提出的 CLEAN\cite{CLEAN2023}（Contrastive Learning-Enabled enzyme ANnotation）基于 ESM-1b 蛋白语言模型与对比学习框架预测酶的 EC 编号，在不依赖同源搜索的条件下达到与 BLAST 相当甚至更高的准确度，是当前基于序列的酶功能预测的代表性领先方法。上述两类方法的共同局限是：EC 编号仅刻画酶的反应类型（催化何种化学变化），并不直接给出酶能否催化某个具体底物，难以支撑精确到分子的特异性预测需求。

\textbf{基于分子图的方法与蛋白--化合物交互建模}。另一类方法将底物的分子结构直接纳入建模，以预测"酶--底物对是否发生催化"的二分类。Tsubaki 等提出的 CPI-prediction\cite{CPI2019} 是该路线的早期代表之一，采用图神经网络对底物分子图编码、卷积网络对蛋白序列编码，两侧特征通过端到端融合预测交互；该方法思想简单而通用，其后的诸多 CPI 与酶--底物预测模型均在此基础上演进。Kroll 等提出的 ESP\cite{ESP2023}（Enzyme Substrate Prediction）基于 ESM-1b 蛋白语言模型与 GNN 分子嵌入组合，在跨物种、跨反应的广谱酶--底物数据上取得显著进展，是当前通用酶--底物二分类任务的重要基线。与 DeepEC、CLEAN 相比，ESP 的粒度从 EC 类转向具体分子，更贴近本研究的任务形式。

\textbf{基于结构的方法与对接复合物建模}。将酶的三维结构纳入预测是最近一代方法的核心方向。AlphaFold2\cite{AlphaFold2} 对未解析结构的蛋白给出近晶体质量的预测构象，使得基于结构的方法的数据瓶颈得到关键突破；结合分子对接（如 AutoDock Vina\cite{Vina2010,Vina2021}、Uni-Dock\cite{UniDock2023}）生成的酶--底物结合姿态，使模型可以直接从三维微环境中学习特异性规律。EZSpecificity\cite{EZSpec2025} 正是这一路线当前最先进的工作（详见 \ref{sec:ezspec-overview} 节），其 ESIBank 数据集与 SE(3)-等变结构通道代表了数据层与模型层两个维度的共同进步。在 P450 特定领域之外，Liu 等在 Nature Catalysis 发表的 EnzymeCAGE\cite{EnzymeCAGE2026} 引入残基级几何向量感知通路，对酶--底物检索任务给出了进一步的结构感知方案，也是本研究第五章双图架构改进的参考实现来源。

\textbf{其他任务形式}。除二分类之外，酶反应动力学参数预测（如 Li 等的 DLKcat\cite{DLKcat2022} 预测 $k_{\mathrm{cat}}$）是另一条互补路线，可为二分类模型提供额外的定量信号，但与本研究的二分类目标正交，在此不展开综述。

综合以上工作，酶底物特异性预测的研究趋势可概括为：从基于序列转向基于三维结构、从 EC 类别预测转向具体分子级别预测、从通用跨家族转向领域专属评估。本研究正是在上述趋势下、以 EZSpecificity 为起点、以 P450 家族为专属评估场景展开。

\section{EZSpecificity 参考模型框架概述}
\label{sec:ezspec-overview}

EZSpecificity\cite{EZSpec2025} 定义的任务形式为：输入酶的氨基酸序列与底物的 SMILES 表达式，输出该酶是否催化该底物的二分类判定。模型结构由三通道并行的特征提取与双向交叉注意力融合两部分构成（详细结构在本章 \ref{sec:representation} 节与第五章 \ref{sec:ch5-design} 节中进一步展开）。

\textbf{三通道并行编码}。序列通道以 ESM-2\cite{ESM2023} 对酶氨基酸序列生成残基级 1{,}280 维嵌入，经线性层压缩至 128 维并做归一化；分子通道由 GROVER\cite{GROVER2020} 的原子级与分子级嵌入与 Morgan 指纹\cite{ECFP2010}（1{,}024 位）组合构成，各自经线性层归一至 128 维；结构通道对酶--底物对接复合物 PDB 中以配体重原子为中心 $10~\mathrm{\AA}$ 范围内的口袋残基（不含氢）与配体原子构图，经 3 层 E($n$)-等变图神经网络（EGNN）\cite{EGNN2021} 消息传递后得到每原子 128 维特征。三通道的核心设计原则是\textbf{分而治之}：每通道在各自的表征空间内独立学习，融合环节留给交叉注意力。

\textbf{双向交叉注意力与预测头}。基于 Vaswani 等提出的多头注意力机制\cite{Transformer2017}（其思想可追溯至 Bahdanau 等最早在神经机器翻译中引入的注意力机制\cite{Bahdanau2015}），EZSpecificity 以 8 头交叉注意力同时执行两个方向：酶残基向底物原子查询，与底物原子向酶残基查询。两方向的注意力输出与三通道原始特征拼接为 7 个 128 维向量（共 896 维），送入多层感知机预测头输出二分类 logits。

\textbf{选定该参考框架的原因}。本研究选择 EZSpecificity 作为参考框架有三方面理由：第一，它是当前基于三维结构的酶--底物特异性预测中最具代表性的领先方法之一，代表现有技术路线的性能上限；第二，其数据管线与特征生成流程完整开源，可直接在 P450 数据集上复现；第三，其三通道架构模块化清晰，便于在其上开展结构通道的方向感知改进（即本研究第五章双图架构）。

\section{蛋白与分子表征学习组件}
\label{sec:representation}

本节集中介绍构成 EZSpecificity 框架的四类表征学习组件，以便读者在不回顾原论文的情况下理解本研究的技术基础。

\textbf{蛋白语言模型}。ESM-2\cite{ESM2023} 是 Lin 等在 Science 2023 发表的蛋白语言模型，以 UniRef90/UniRef50（约 6{,}500 万条 UniRef50 簇代表序列）的蛋白序列为训练语料，以掩码语言建模目标预训练。相较于第一代 ESM\cite{ESM2021}，ESM-2 在模型规模、训练数据与下游任务（包括基于 ESMFold 的端到端结构预测）支撑能力上均有显著提升，输出的残基级嵌入已被证明能够隐式编码蛋白的二级结构、折叠家族与功能位点信息。本研究与 EZSpecificity 参考框架均采用 ESM-2 作为酶序列通道的特征提取器。

\textbf{分子指纹与分子图表征}。Morgan 扩展连接指纹\cite{ECFP2010} 在 Morgan 1965 年提出的原子环境迭代算法基础上，由 Rogers 与 Hahn 于 2010 年扩展为 ECFP（Extended-Connectivity Fingerprints），以化学键半径逐步展开编码分子的局部子结构，是药物化学领域最广泛使用的分子特征之一；作为规则式表征，它在小数据场景下仍具备稳定的基线性能。GROVER\cite{GROVER2020} 是 Rong 等在 NeurIPS 2020 提出的分子图变换器，在千万级未标注分子上采用自监督预训练，以双层次 Transformer 分别建模原子级与分子级表征；其原子级与分子级嵌入为下游任务提供互补的学习型分子特征。本研究与 EZSpecificity 将 Morgan 指纹与 GROVER 嵌入组合使用以兼顾规则特征与学习特征。

\textbf{几何表征与等变图神经网络}。处理蛋白--分子三维结构需要满足两类空间对称性：旋转--平移不变（SE(3)-invariant），即整体刚体变换不改变模型输出；以及旋转--平移等变（SE(3)-equivariant），即变换下输出以相应方式协同变化。EGNN\cite{EGNN2021} 是 Satorras 等提出的 E($n$)-等变 GNN 层（坐标对旋转--平移--反射等变、标量特征对其不变），以原子间距离为核心驱动消息传递，在保持对整体刚体变换鲁棒的同时避免了高阶等变表征带来的计算开销；EZSpecificity 与本研究基线模型均以 EGNN 构成结构通道。GVP\cite{GVP2021,GVPMacro2021} 是 Jing 等提出的 SE(3)-等变层，区别于 EGNN 的关键在于同时处理标量与向量两类特征：标量部分按常规神经网络变换，向量部分在整体旋转下跟随旋转，从而保留了 EGNN 所抛弃的方向信息。本研究第五章的双图架构改进正是在 EGNN 原子级通路之上并行引入 GVP 残基级通路，以尝试补足原架构的方向感知缺失。

\textbf{注意力机制}。注意力机制最早由 Bahdanau 等在神经机器翻译中提出\cite{Bahdanau2015}，以软对齐替代传统的固定上下文向量；Vaswani 等 2017 年在此基础上提出的 Transformer\cite{Transformer2017} 将多头自注意力与前馈网络组合成纯注意力架构，成为现代深度学习的通用基础。EZSpecificity 的双向交叉注意力机制在两个方向上相互应用多头注意力，实现酶残基与底物原子的互相查询。本研究第五章在此基础上叠加残基级 GVP 分支，不改变交叉注意力本身的结构。

\section{研究问题与总体方案}
\label{sec:problem-formulation}

\textbf{任务形式化}。本研究采用与 EZSpecificity 一致的任务定义：给定酶 $e$（以氨基酸序列表示）与底物 $s$（以 SMILES 表示）及其对接复合物结构 $c(e, s)$，构造二分类函数 $f(e, s, c) \in \{0, 1\}$，输出该酶是否催化该底物的判定。本研究以 P450 为目标家族，构造领域专属的训练、验证与测试数据集，并在测试集上评估模型性能。

\textbf{评价指标}。鉴于数据集的正负样本不均衡（正负比约 $1{:}10$），本研究采用 AUC-ROC 与 AUPR 两项指标并列评估：AUC-ROC 在成对排序意义下评估模型的总体区分能力，但 Saito 与 Rehmsmeier\cite{Saito2015} 指出，在类别极端不平衡时 AUC-ROC 的绝对数值可能偏乐观、单独使用易产生误导；AUPR 在稀有正样本打分分布上更敏感，其随机基线即正样本占比，能更直接度量模型在少数类上的真实表现。两指标并列报告以避免单一指标造成的评估偏差。

\textbf{外部公平对比协议}。为评估 EZSpecificity 公开 checkpoint 在 P450 上对未见酶的泛化能力，本研究在测试集基础上额外构造过滤后外部测试集：以 InterPro\cite{Blum2025} 功能域严格识别 EZSpecificity 原始训练集 ESIBank\cite{EZSpec2025} 中的 389 个真正 P450，将与之酶层面重合的测试样本剔除。过滤后得到的 7{,}963 条测试集将参考模型的评估条件收敛为"完全未见过的 P450 酶"，为泄漏控制的公平对比提供数据基础。

\textbf{总体研究方案}。本研究按三步推进：（1）第三章构建多源整合的 P450 专属基准数据集，并生成与 EZSpecificity 一致的特征管线；（2）第四章在该基准上从头训练同架构基线模型 M0，并与 EZSpecificity 公开 checkpoint 在过滤后测试集上进行泄漏控制公平对比；（3）第五章在原架构基础上引入残基级 GVP 通路构成双图架构 M1，在单 random fold 评估下作为方向信息感知设计的初步可行性探索。全部训练与评估严格限定在 random split 第一折上完成，不开展交叉验证或多种子稳健性重复；相应局限在第六章如实声明。

\textbf{与相关工作的区别}。本研究在综述的相关工作基础上形成如下差异：在数据层面，明确聚焦于 P450 家族并整合既往 ML 基准中较少覆盖的植物 P450 来源；在评估协议层面，引入 InterPro 严格过滤的泄漏控制外部测试集，为 SOTA 通用模型在稀缺家族上的真实能力提供公平度量；在架构层面，针对 P450 催化对方向信息敏感的物理特性，尝试在原子级 EGNN 通路之外引入残基级 GVP 通路。相关工作中述及的 CLEAN、ESP、DeepEC、CPI-prediction、DLKcat 等方法因任务目标、数据规模或评估协议的差异，未作为本研究的直接基线；相应的方法学定位与对比在第四章与第六章中进一步讨论。
